\documentclass{exam}

\newif\ifA
\newif\ifkey
\newif\ifprintselected   % required by exam_config's \ansbox

%%%%%%%%%%%%%%%%%%%%%%%%%%%%%%%%%%%%%%%%%%%%%%%%%%%%%%
%%% SET THESE IF VARIABLES
%%%%%%%%%%%%%%%%%%%%%%%%%%%%%%%%%%%%%%%%%%%%%%%%%%%%%%
% Lab name -- labs are named by the lesson they belong to, never by a
% sequential lab number (Neil, 2026-10-12): this is the L18--L19 Lab.
\def\labLesson{18--L19}

% is this the key or not?
\ifdefined\iskey
	\ifnum\iskey=1
		\keytrue
	\else
		\keyfalse
	\fi
\else
	%% Pick one manually
	\keyfalse % if this is NOT the key
%	\keytrue     % if this IS the key
\fi

% set up the header and footer - change for every term
\ifdefined\thisTermIn
	\def\thisterm{\thisTermIn}
\else
	\def\thisterm{Spring 2027}
\fi
%
\def\thisclass{ECE 448}

%%%%%%%%%%%%%%%%%%%%%%%%%%%%%%%%%%%%%%%%%%%%%%%%%%%%%%
%%% END SET THESE IF VARIABLES
%%%%%%%%%%%%%%%%%%%%%%%%%%%%%%%%%%%%%%%%%%%%%%%%%%%%%%
% packages and shortcuts
\input{myPackages_gr}
\usepackage[hidelinks, linkcolor=red]{hyperref}
\input{myShortcuts}
\setlength{\parindent}{0pt}
\setlength{\parskip}{6pt}

%% exam class customization
\input{exam_config}
% ECE 448 document fonts.
%
% Body text is Barlow (SIL OFL; the faces and the license are in
% latex/fonts/Barlow/). Barlow is not in TeX Live, so it is loaded as an
% OpenType face through fontspec -- which means the document must be built
% with lualatex, not pdflatex. build_practice.sh and build_lab.sh both do.
%
% Compiling with pdflatex anyway is not an error: the guard below falls back
% to the LaTeX default so the source still compiles under any engine and can
% be syntax-checked. The resulting PDF just is not the one we ship.
%
% Paths are relative to latex/, which is where both build scripts run
% (each does cd "$(dirname "$0")"). Loading by path rather than by installed
% family name means no system font install and no fontconfig, so a fresh
% clone builds a byte-comparable PDF anywhere.
%
% MATH IS DELIBERATELY LEFT IN COMPUTER MODERN. Barlow has no math companion,
% and CM math under sans body text is already the house pairing: the decks set
% words in Source Sans Pro and let MathJax render symbols in its CM-derived
% face. Moving math to a sans as well would mean unicode-math plus a sans math
% font (Fira Math being the plausible one) and re-checking every equation in
% the course for symbol coverage. That is a separate decision, not a side
% effect of this one.

\usepackage{iftex}
\ifLuaTeX
	\usepackage{fontspec}
	\defaultfontfeatures{Ligatures=TeX}
	\setmainfont{Barlow}[
		Path           = fonts/Barlow/,
		Extension      = .ttf,
		UprightFont    = *-Regular,
		ItalicFont     = *-Italic,
		BoldFont       = *-Bold,
		BoldItalicFont = *-BoldItalic,
	]
	\setsansfont{Barlow}[
		Path           = fonts/Barlow/,
		Extension      = .ttf,
		UprightFont    = *-Regular,
		ItalicFont     = *-Italic,
		BoldFont       = *-SemiBold,
		BoldItalicFont = *-BoldItalic,
	]
	% microtype earns its place here specifically because of Barlow. Its
	% glyphs are narrower than Computer Modern's, so justified prose sets
	% tighter and TeX has to squeeze interword space to fit; two paragraphs
	% in the L02/L03 sets went overfull on the first Barlow build. Protrusion
	% and font expansion absorb that -- both are lualatex features, so this
	% belongs inside the guard with everything else.
	\usepackage{microtype}

	% Mono stays Latin Modern: Barlow ships no monospaced face, and verbatim
	% and code need one that actually is fixed-pitch.
\fi
   % Barlow body text (needs lualatex)
\printselectedfalse      % full worked solutions in the key, not answers-only
\CorrectChoiceEmphasis{\color{red}}
\SolutionEmphasis{\color{red}}
\renewcommand\questionlabel{\textbf{\thequestion.}}
\renewcommand{\thepartno}{\alph{partno}}
\renewcommand{\partlabel}{(\thepartno)}

% Fillable table cell: blank in the student copy, the expected value in red on
% the KEY. Fixed-width centered columns keep the blank cells writable.
\newcolumntype{K}[1]{>{\centering\arraybackslash}p{#1}}
\newcommand{\kv}[1]{\ifkey{\color{red}#1}\fi\rule[-0.10in]{0pt}{0.34in}}

% Answers are not in this public file. Every solution, \kv value, and \ansbox
% value is \keysol{<id>}, defined in the private ECE448-faculty repo at
% lab-keys/src/ECE448_Lab_L18_ADSBAIS_solutions.tex; build_lab.sh puts that
% folder on TEXINPUTS (KEYDIR) for the KEY. A KEY build without it stops.
\long\def\defkeysol#1#2{\expandafter\long\expandafter\def\csname keysol@#1\endcsname{#2}}
\def\keysol#1{\ifcsname keysol@#1\endcsname\csname keysol@#1\endcsname\fi}
\ifkey
	\IfFileExists{ECE448_Lab_L18_ADSBAIS_solutions.tex}%
		{\input{ECE448_Lab_L18_ADSBAIS_solutions}}%
		{\errmessage{KEY build needs ECE448_Lab_L18_ADSBAIS_solutions.tex from ECE448-faculty/lab-keys/src (set KEYDIR)}}
\fi

% print the answers if this is the key
\ifkey
	\printanswers % if this IS the key
\else
	\noprintanswers  %if this is NOT the key
\fi

% print the header and footer
% need to print key on top if this is the key
\ifkey
	\firstpageheader{}{\color{red}{ KEY}}{}
	\runningheader{\thisclass\hspace{1pt} -- \thisterm }{L\labLesson\ Lab\hspace{1pt}\color{red}{ KEY}}{\thepage\hspace{1pt} of\hspace{1pt} \numpages}
	\firstpagefooter{}{}{}
	\runningfooter{}{}{}
\else
	\firstpageheader{}{}{}
	\runningheader{\thisclass \hspace{1pt} -- \thisterm }{L\labLesson\ Lab\hspace{1pt}}{\thepage\hspace{1pt} of\hspace{1pt} \numpages}
	\firstpagefooter{}{}{}
	\runningfooter{}{}{}
\fi

\runningheadrule
\begin{document}
\pagestyle{headandfoot}   % exam class

%%%%%%%%%%%%%%%%%%%%%%%%%%%%%%%%%%%%%%%%%%%%%%%%%%%%%%
% title block
%%%%%%%%%%%%%%%%%%%%%%%%%%%%%%%%%%%%%%%%%%%%%%%%%%%%%%

\begin{center}
USAF ACADEMY \\
DEPARTMENT OF ELECTRICAL AND COMPUTER ENGINEERING \\
LESSONS 18--19 LAB - ADS-B AND AIS \\
\textbf{Lab Packet} \\
\thisterm \\[4mm]
\end{center}

\vspace{-4pt}
\textbf{Name:} \rule[-2pt]{2.6in}{0.4pt} \hfill
\textbf{Section:} \rule[-2pt]{0.9in}{0.4pt} \hfill
\textbf{Date:} \rule[-2pt]{1.1in}{0.4pt}

\vspace{4pt}

\fullwidth{\textbf{Learning Objective 6.7: } I can receive and decode ADS-B messages from
aircraft transponders.\\
\textbf{Learning Objective 6.8: } I can receive and decode AIS messages from marine vessels.}

\vspace{2pt}

This packet is the L18--L19 Lab, in three parts, and you turn it in once, complete,
before Lesson 20. Do \textbf{Part 1} on your own before the lab period: fly an aircraft of your
own through the ADS-B generator and work out the settings for the lab's AIS recording. \textbf{Part 2} is your
own measurements, recorded in class: Questions 3 to 5 with a PlutoSDR, gr-adsb, and live aircraft,
and Questions 6 to 8 on the lab's AIS recording. Do \textbf{Part 3} on your own after the
lab: explain what you saw.

\textbf{Grading.} You submit the packet once, we grade it against Objectives 6.7 and 6.8
on the course's 1 / 0 scale, and you may revise it under course policy. Paste screenshots
where a question asks for one.

\textbf{Never on 1090 MHz.} The Pluto only receives on 1090 MHz, and the AIS recording is
simulated, written by a GNU Radio flowgraph rather than captured from an antenna. Never
transmit the ADS-B or AIS generators on 1090 MHz or on the AIS channels: they carry
safety-of-life traffic. The one transmission in this lab is optional, in Question 4: your
own aircraft from \texttt{L18-adsb-tx-pluto.grc}, sent at 915 MHz into a cable from the
Pluto's TX port through a 30 dB attenuator, or, if your instructor says so, over the air at
915 MHz at low power. Never set TX attenuation below 30 dB.

The flowgraphs are on the L18 and L19 lesson pages: \texttt{L18-adsb-tx.grc},
\texttt{L18-adsb-tx-pluto.grc}, \texttt{L18-adsb-rx.grc}, and \texttt{L19-ais-rx.grc}.
\texttt{L18-adsb-rx.grc} needs gr-adsb, which radioconda includes as \texttt{gnuradio-adsb}. The lab's AIS recording,
\texttt{L18-L19-lab-ais.sigmf-data} and its \texttt{.sigmf-meta}, is on the lab page.
Its meta file gives the datatype \texttt{cf32\_le}, a sample rate of 192000, and a
center frequency of 162025000.

\begin{questions}

\fullwidth{\large\textbf{Part 1: Pre-Lab} \normalsize(individual, before the lab period)}

%%%%%%%%%%%%%%%%%%%%%%%%%%%%%%%%%%%%%%%%%%%%%%%%%%%
\question \textbf{Your own aircraft.} Put \texttt{L18-adsb-tx.grc} and
\texttt{L18-adsb-rx.grc} in one folder. Double-click the generator's ADS-B Out block,
choose Open in Editor, and add a fifth line to \texttt{AIRCRAFT}: an ICAO address from
\texttt{0xA00000} to \texttt{0xAFFFFF} that you choose, a callsign of up to 8 capital
letters and digits made from your initials, a position within 50~km of USAFA, and an
altitude, speed, track, and climb rate of your choice, and an amplitude between 0.45
and 0.8 that is not 0.6, so that no aircraft in the list has it. Save, run the generator once, then run the receiver.

\begin{parts}

	\part Write the line you added, and the row the receiver's table shows for your
	aircraft. Do they agree?
		\begin{solution}[1.0in]\keysol{sol1}\end{solution}

	\part Raise the Added noise slider a step at a time until your aircraft leaves the
	table, and record that noise level. Paste a screenshot of the receiver just before it
	leaves.
		\begin{solution}[1.5in]\keysol{sol2}\end{solution}

\end{parts}

%%%%%%%%%%%%%%%%%%%%%%%%%%%%%%%%%%%%%%%%%%%%%%%%%%%
\newpage
\question \textbf{AIS settings for the lab recording.} Your AIS receiver in Question 6
plays \texttt{L18-L19-lab-ais.sigmf-data}, whose meta file is described above.

\begin{parts}

	\begin{minipage}[t]{0.58\linewidth}
		\part Find the center frequency of each channel's Frequency Xlating FIR Filter.
		\begin{solution}[0.9in]\keysol{sol3}\end{solution}
	\end{minipage}\hfill%
	\begin{minipage}[t]{0.37\linewidth}
		\raggedleft \vspace{12pt}
		A: \ansbox[1.0in]{\keysol{ans1}}

		\vspace{4pt}
		B: \ansbox[1.0in]{\keysol{ans2}}
	\end{minipage}

	\begin{minipage}[t]{0.58\linewidth}
		\part Choose \texttt{decim} for 48 kS/s after the filter, and find the samples per
		bit and the Quadrature Demod gain.
		\begin{solution}[1.0in]\keysol{sol4}\end{solution}
	\end{minipage}\hfill%
	\begin{minipage}[t]{0.37\linewidth}
		\raggedleft \vspace{12pt}
		decim: \ansbox[0.9in]{\keysol{ans3}}

		\vspace{4pt}
		sps: \ansbox[0.9in]{\keysol{ans4}}

		\vspace{4pt}
		gain: \ansbox[0.9in]{\keysol{ans5}}
	\end{minipage}

\end{parts}

%%%%%%%%%%%%%%%%%%%%%%%%%%%%%%%%%%%%%%%%%%%%%%%%%%%
\newpage
\fullwidth{\large\textbf{Part 2: Measurements} \normalsize(in class)}

\question \textbf{Build an ADS-B front end.} Start GRC with an empty canvas and build the
front end the lab page lists: a PlutoSDR Source at 1090 MHz and 2 MS/s, an AGC3, Complex
to Mag\^{}2, and a triggered QT GUI Time Sink.

\begin{parts}

	\part Paste a screenshot of your flowgraph on the GRC canvas. Every connection must be
	black, not red.
		\begin{solution}[1.4in]\keysol{sol5}\end{solution}

	\part Paste a screenshot of the time sink frozen on a live preamble. Record the trigger
	level you chose, the noise floor, and the pulse heights of the strongest and weakest
	messages you saw, all in power, as the plot shows them.
		\begin{solution}[1.6in]\keysol{sol6}\end{solution}

\end{parts}

%%%%%%%%%%%%%%%%%%%%%%%%%%%%%%%%%%%%%%%%%%%%%%%%%%%
\newpage
\question \textbf{Live aircraft.} Add gr-adsb's ADS-B Framer, Demod, and Decoder after
Complex to Mag\^{}2, copy the Aircraft Table and its two Function Probes and QT GUI Labels from
\texttt{L18-adsb-rx.grc}, and run it for 10 minutes with the Pluto at a window, as the lab
page lists. If no aircraft decode, decode your own Question 1 aircraft sent from
\texttt{L18-adsb-tx-pluto.grc} at 915 MHz into a cable through the 30 dB attenuator (see
the box above), or the lesson's recording, and say which.

\begin{parts}

	\part Record up to five aircraft, and find each one on a public flight tracker such as
	ADS-B Exchange or FlightAware. Mark whether the tracker agrees. For your own aircraft or
	the recording, compare with the \texttt{AIRCRAFT} list instead.

	\begin{center}
	\renewcommand{\arraystretch}{1.3}
	\small
	\begin{tabular}{| K{0.65in} | K{0.8in} | K{0.7in} | K{0.6in} | K{1.2in} | K{0.7in} |}
	\hline
	ICAO & Callsign & Altitude (ft) & Speed (kt) & Latitude, longitude & Tracker agrees? \\\hline
	& & & & & \\\hline
	& & & & & \\\hline
	& & & & & \\\hline
	& & & & & \\\hline
	& & & & & \\\hline
	\end{tabular}
	\end{center}
		\begin{solution}[0.6in]\keysol{sol7}\end{solution}

	\part Paste a screenshot of the running window with the Messages label and the table.
		\begin{solution}[1.4in]\keysol{sol8}\end{solution}

\end{parts}

%%%%%%%%%%%%%%%%%%%%%%%%%%%%%%%%%%%%%%%%%%%%%%%%%%%
\newpage
\question \textbf{Gain.} With the decoder running live, set the Pluto RX gain to each value
in turn, wait 60~s, and record the number decoded on the Messages label (messages the
Decoder turned into an identity, a position, or a velocity) and the number of aircraft in
the table.

\begin{center}
\renewcommand{\arraystretch}{1.3}
\begin{tabular}{| l | K{0.8in} | K{0.8in} | K{0.8in} | K{0.8in} | K{0.8in} |}
\hline
Pluto RX gain (dB) & 70 & 60 & 50 & 40 & 30 \\\hline
Decoded in 60 s & & & & & \\\hline
Aircraft & & & & & \\\hline
\end{tabular}
\end{center}

\begin{parts}
	\part Describe which aircraft dropped out first as the gain fell, and how their
	altitude or position compared with the ones that stayed.
		\begin{solution}[1.3in]\keysol{sol9}\end{solution}
\end{parts}

%%%%%%%%%%%%%%%%%%%%%%%%%%%%%%%%%%%%%%%%%%%%%%%%%%%
\newpage
\question \textbf{Build an AIS receiver.} Start GRC with an empty canvas and build the
two-channel receiver the lab page lists for \texttt{L18-L19-lab-ais.sigmf-data}, with
your Question 2 settings and the Squelch slider at $-30$ dB.

\begin{parts}

	\part Paste a screenshot of your flowgraph on the GRC canvas. Every connection must be
	black, not red.
		\begin{solution}[1.5in]\keysol{sol10}\end{solution}

	\begin{minipage}[t]{0.58\linewidth}
		\part Run it for one pass through the 5~s recording and record the Frames label.
		How many ships are in the table?
		\begin{solution}[0.9in]\keysol{sol11}\end{solution}
	\end{minipage}\hfill%
	\begin{minipage}[t]{0.37\linewidth}
		\raggedleft \vspace{12pt}
		ships: \ansbox[0.9in]{\keysol{ans6}}
	\end{minipage}

\end{parts}

%%%%%%%%%%%%%%%%%%%%%%%%%%%%%%%%%%%%%%%%%%%%%%%%%%%
\newpage
\question \textbf{The ships.} Read the decoder's printout in the GRC console and the table.

\begin{parts}

	\part Record every ship. The flag country comes from the first three digits of the
	MMSI, its Maritime Identification Digits; look them up in the ITU's MID table.

	\begin{center}
	\renewcommand{\arraystretch}{1.3}
	\footnotesize\setlength{\tabcolsep}{3pt}
	\begin{tabular}{| K{0.75in} | K{0.95in} | K{0.75in} | K{0.75in} | K{1.1in} | K{0.75in} |}
	\hline
	MMSI & Name & Flag & Destination & Latitude, longitude & Speed, course \\\hline
	\kv{\keysol{sh1a}} & \kv{\keysol{sh1b}} & \kv{\keysol{sh1c}} & \kv{\keysol{sh1d}} & \kv{\keysol{sh1e}} & \kv{\keysol{sh1f}} \\\hline
	\kv{\keysol{sh2a}} & \kv{\keysol{sh2b}} & \kv{\keysol{sh2c}} & \kv{\keysol{sh2d}} & \kv{\keysol{sh2e}} & \kv{\keysol{sh2f}} \\\hline
	\kv{\keysol{sh3a}} & \kv{\keysol{sh3b}} & \kv{\keysol{sh3c}} & \kv{\keysol{sh3d}} & \kv{\keysol{sh3e}} & \kv{\keysol{sh3f}} \\\hline
	\kv{\keysol{sh4a}} & \kv{\keysol{sh4b}} & \kv{\keysol{sh4c}} & \kv{\keysol{sh4d}} & \kv{\keysol{sh4e}} & \kv{\keysol{sh4f}} \\\hline
	\end{tabular}
	\end{center}
		\begin{solution}[0.5in]\keysol{sol12}\end{solution}

	\part Copy one complete \texttt{!AIVDM} sentence from the console, paste it into
	\texttt{pyais} or an online AIVDM decoder, and record whether its fields match your
	table.
		\begin{solution}[1.0in]\keysol{sol13}\end{solution}

\end{parts}

%%%%%%%%%%%%%%%%%%%%%%%%%%%%%%%%%%%%%%%%%%%%%%%%%%%
\newpage
\question \textbf{The squelch.} Set the Squelch slider to each value in turn, wait for
two passes through the recording, and mark which ships are in the table.

\begin{center}
\renewcommand{\arraystretch}{1.25}
\small\setlength{\tabcolsep}{3pt}
\begin{tabular}{| l | *{7}{K{0.62in} |}}
\hline
Squelch (dB) & $-40$ & $-30$ & $-20$ & $-15$ & $-10$ & $-5$ & 0 \\\hline
Ships decoded & \kv{\keysol{sq1}} & \kv{\keysol{sq2}} & \kv{\keysol{sq3}} & \kv{\keysol{sq4}} & \kv{\keysol{sq5}} & \kv{\keysol{sq6}} & \kv{\keysol{sq7}} \\\hline
\end{tabular}
\end{center}

\begin{parts}
	\part Record the Frames label after one pass at $-40$ dB and at $-30$ dB, and describe
	the time plot at each.
		\begin{solution}[1.2in]\keysol{sol14}\end{solution}
\end{parts}

%%%%%%%%%%%%%%%%%%%%%%%%%%%%%%%%%%%%%%%%%%%%%%%%%%%
\newpage
\fullwidth{\large\textbf{Part 3: Writeup} \normalsize(individual, after the lab)}

\question \textbf{Analysis.}

\begin{parts}

	\part Compare your two receivers, using what your own time sinks showed. Explain why
	the ADS-B receiver needs only the power of each sample, while the AIS receiver needs
	an FM demodulator and timing recovery.
		\begin{solution}[1.5in]\keysol{sol15}\end{solution}

	\part Use your squelch table to rank the ships you decoded by received strength, and explain
	what each step of the table tells you. Then explain how you would choose the squelch
	for a real recording.
		\begin{solution}[1.5in]\keysol{sol16}\end{solution}

	\part Compare your live aircraft with the flight tracker, row by row: for each, give
	the difference in altitude and in position you saw, and explain it. If you decoded your
	own aircraft or the recording, compare with the \texttt{AIRCRAFT} list instead.
		\begin{solution}[1.5in]\keysol{sol17}\end{solution}

\end{parts}

\end{questions}

\vspace{1.0em}
\noindent\textbf{Documentation:} \rule[-2pt]{0.6\linewidth}{0.4pt}

\end{document}
